% Selected-scale charts and the public gallery handoff.
\section{From funded chart receipts to a trace gallery}
\label{sec:ex-lattice-firesky-gallery}

This section records the actual selected-run route from one-scale chart data
to the scale-dependent charts used by the trace calibration. The analytic
chart receipts are provided by the funded construction and are identified
below. We prove the subsequent scale selection and finite-dimensional limit
steps explicitly. The correction and Ward-flux calculations are in
Section~\ref{sec:ex-lattice-trace-producer}.

\subsection{The funded chart family and its one-step output}

Fix a grade $d$, a binding, a gauge-invariant chart budget, and a scale
parameter $z\geq0$. The funded chart receipt consists of three panels
$P(t,i)$, $i=0,1,2$, for $0\leq t\leq1$, with arc coordinates $b_i+t$.
Their currency and angular parameters are constant in $t$, and the three
base arcs obey the finite separation inequalities for the probe count at
grade $d$. At every sufficiently large level $k$ and every $t\in[0,1]$,
the receipt supplies a finite-volume chart witness at that level and
offset. The other fields of this witness are the chart, invoice,
legacy-cost, actuation, and stamp data used by the quantitative chart
estimate. The single offset changes only the three arc coordinates.

The simultaneous choice needed here is stronger than choosing a funded
run separately at every grade. We give its construction next. The local
analytic inputs are stated separately, so that the choice of a common
schedule does not hide an estimate.

\begin{lemma}[A quantitative, signed-equivariant reading]
\label{lem:trace-quantitative-invoice}
Keep the funded hypotheses and notation of
Proposition~\ref{prop:rank-selected-tested-readout}. At each fixed grade
$d$, its reading can be chosen so that, for all sufficiently large levels,
uniformly in $N,p,u$,
\begin{equation}\label{eq:trace-quantitative-invoice}
 |Au-Tu|_1\le H^{-1/2}|Au|_1,\qquad H=\mathcal M_k.
\end{equation}
Assume that the calibrated map $A$ intertwines the signed coordinate
symmetries of its label space. The reading can then also be chosen to
intertwine those symmetries, preserving its exact unit and real structure.
\end{lemma}
\begin{proof}
The proof of Lemma~\ref{lem:rank-growing-wick-pin} gives, before discarding
its rate,
\[
 |Au-Tu|_1\le c_d\{H^{-1}+D2^DH^{-3/4}
                 +w_DH^{-1+\alpha(D+1)}\}|Au|_1,
 \qquad \alpha=\frac1{4(D+2)}.
\]
Since $-1+\alpha(D+1)<-3/4$, this is at most
$C_dH^{-3/4}|Au|_1$, with
$C_d=c_d(1+D2^D+w_D)$, for $H\ge1$. Thus it is at most
$H^{-1/2}|Au|_1$ once $H^{1/4}\ge C_d$.

For the symmetry assertion, let $\Sigma$ be the finite signed coordinate
group, with representation $\rho$ on the grade-$d$ label space, and put
$a_d=\max(1,\max_{\sigma\in\Sigma}\|\rho(\sigma)\|_{1\to1})$.
Set
\[
 \overline T=\frac1{|\Sigma|}\sum_{\sigma\in\Sigma}
                \rho(\sigma)T\rho(\sigma)^{-1}.
\]
Reindexing this finite sum proves equivariance. Since $A\rho(\sigma)
=\rho(\sigma)A$, each summand's error on $u$ is bounded by
$a_d^2 C_dH^{-3/4}|Au|_1$. The triangle inequality gives the same bound
for the average. Enlarging the common threshold until
$H^{1/4}\ge a_d^2C_d$ proves \eqref{eq:trace-quantitative-invoice}
for $\overline T$. Every summand fixes the unit, and signed coordinate
actions commute with complex conjugation, so these identities survive
the average. The finite group and constants depend only on the grade,
not on $k,N,p$. Hence taking the maximum for $0\le d'\le d$ gives
one finite-cutoff threshold.
\end{proof}

The averaging of a reading must also be performed on its raw and Wick
polynomials when a stamp is required. One cannot change $T$ alone and
retain a stamp for the old reading. Signed pullback preserves the healthy
cube and the halo; it carries the raw polynomial, its Wick polynomial and
their reading together because the atlas and its covariance have the
same signed symmetry. The finite sum commutes with these linear
operations. The degree bound is unchanged, and coefficient and ticket
bounds follow by the triangle inequality. Their fixed symmetry constants
must be included in the local price purchase. The bound above has the
explicit $H^{1/4}$ margin that absorbs those constants in the invoice;
the raw-polynomial and healthy-ticket bounds must retain their own
reserved constants from the selected-family construction.

\paragraph{Local inputs for the simultaneous purchase.}
Write $p_j$ for the order of the stage-$j$ profile and assume $p_j\to\infty$.
For each $d$ choose three base arcs $b_{d,i}\ge d+6$, $i=0,1,2$, satisfying
\eqref{eq:three-checks}, with arity $n_d=\operatorname{trays}(d)\ge2$.
The explicit choice after that equation gives such arcs. Purchase the
three panels initially at the larger arcs $b_{d,i}+1$. The local panel
purchase has the following precise interface.
\begin{enumerate}
\item The panels have one common Taylor order $\vartheta_d\ge d$ and,
 at each stage $j$, one common polynomial reading order $Q_{d,j}$.
 Their remaining constants, including currencies, are fixed before
 choosing a scale schedule. If $\vartheta_d+1\le p_{j-1}$, there is a
 block floor $\mu_{d,i,j}$ such that the local comparison estimates
 apply whenever both adjacent block lengths $L_{j-1},L_j$ exceed
 $\mu_{d,i,j}$. The Coulomb estimates for the same panels also hold.
\item After those lengths have been fixed, one may decrease the positive
 local accuracy parameter of each panel, preserving its arc, Taylor
 order and reading order, so that all its local capacity estimates
 hold for every sufficiently large coupling. This threshold is chosen
 at fixed $(d,i,j)$, before the background history and its finite volumes.
 The estimates are uniform over every reduced arc in
 $[0,b_{d,i}+1]$. They supply the finite-volume legacy costs whenever
 the chart window has the dilation, depth and bracket of (F1), with
 coupling between $3\widehat\beta_k/4$ and $\widehat\beta_k$ throughout
 that window. They include the adjacent-stage estimates for each pair
 of stages in $\{j-1,j\}$ that can occur on that window.
\item On any one resulting background history, at fixed $d$ and common
 Taylor order $\vartheta_d$, the local chart and stamp construction is
 valid at every sufficiently large level. Its threshold is uniform
 over the three panels and over reduced arcs. The stamp consists of its raw polynomial, its Wick polynomial and
 its linear reading, with degree and coefficient bounds and a local
 ticket. The representation enters this ticket only through its budget
 radius and healthy threshold; its halo is $Q_{d,j}$ times that radius.
 The covariance and the paid label norm are those of the one selected
 history. These data can therefore be reused in all three panels, as
 checked below; independently chosen readings need not coincide.
\end{enumerate}
These are local analytic hypotheses, not consequences of the
scale-selection argument. In particular item 2 requires the component
comparison estimates. The half-accuracy and stamp portion of item 3 is
supplied by Proposition~\ref{prop:rank-selected-tested-readout}, under
its funded calibrated-family hypotheses. The quantitative invoice is supplied by
Lemma~\ref{lem:trace-quantitative-invoice}; the actuation estimates and
the reserved constants for the signed stamp remain local inputs here.
Mere existence of three invertible maps does not imply them.

\begin{proposition}[One request for all grades and offsets]
\label{prop:trace-joint-chart-purchase}
Assume the preceding local inputs on the common profiles, together with
the common paid scale construction of
Section~\ref{sec:common-scale-selection}. There is one enlargement of its
block floors, entry thresholds and sublinear comparison windows such that
every history on the resulting common calendar has the following property.
Here paying the requirements includes the constant-stage comparison
intervals of \eqref{eq:common-long-plateau}, not merely the block and
entry inequalities.
For every fixed $d$, at every sufficiently large level $k$, the three
panels are available simultaneously at arcs $b_{d,i}+t$ for every
$t\in[0,1]$, with fixed currency and Taylor order. All their charts use
one reading at each grade through $d$. The background history and its
requirements are chosen before any prescribed null sequence of radii.
\end{proposition}
\begin{proof}
\emph{Active grades and block lengths.}
Set, with the predecessor $0-1$ interpreted as $0$,
\[
 K_j=\{d\le j:\vartheta_d+1\le p_{j-1}\},\qquad
 M_j=\max\bigl(\{0\}\cup
       \{\mu_{d,i,j}:d\in K_j,\ 0\le i\le2\}\bigr).
\]
Each $K_j$ is finite. For every fixed $d$, the convergence $p_j\to\infty$
shows that $d\in K_j$ for all sufficiently large $j$; no monotonicity of
$p_j$ is needed for this assertion. Enlarge the block floor at stage $j$
by $\max(M_j,M_{j+1})$, and take a nondecreasing odd sequence above the
enlarged floors, as in the common construction. If $d\in K_j$, then
$L_j\ge M_j\ge\mu_{d,i,j}$ and
$L_{j-1}\ge M_j\ge\mu_{d,i,j}$, including $j=0$ by convention.
Thus one choice of block lengths pays the two adjacent lengths for every
active panel. The maximum with $M_{j+1}$ is essential at a stage change.

\emph{Entry thresholds.}
Now apply local input 2 at these actual lengths. Denote by
$h_{d,i,j}$ a coupling threshold beyond which its conclusion holds for
the resulting, trimmed panel. The trimming leaves $\vartheta_d$ and
$Q_{d,j}$ unchanged, so the common-reading condition is retained.
For inactive triples leave the accuracy parameter unchanged; all these
choices are made before choosing any history. Let
$x(s)=\beta_*+s_Gs$. At a fixed stage $j$, increase the scale-entry
threshold until $x(s)\ge h_{d,i,j+1}$ for every
$d\in K_{j+1}$ and $i=0,1,2$. This is possible because $s_G>0$ and
there are finitely many demands. Include also the finitely many local
capacity and geometric thresholds described below and all the previous
entry thresholds. Taking their maximum preserves every conclusion.

\emph{One comparison window.}
For a fixed block length $L>1$, nonnegative arc $b$, and sufficiently
large $x$, put
\[
 \mathcal M(x)=\exp((\log\log x)^2),\quad
 C(x)=\lceil(\log x)^2\rceil,
 \quad
 h_{L,b}(x)=
 \left\lceil\frac{(b+2)\log\mathcal M(x)}{\log L}\right\rceil
 +\left\lceil\frac{10\log C(x)}{\log L}\right\rceil.
\]
Extend this as any nonnegative integer-valued function below its large-$x$
domain. For fixed $L,b$,
$h_{L,b}(x)=O(1+(\log\log x)^2)$, hence
$h_{L,b}(x(s))/s\to0$. Indeed
$C(x)\le2(\log x)^2$ eventually, so
$\log C(x)\le\log2+2\log\log x$, and rounding adds at most two.
If $w_j^0(s)=o(s)$ is the old stage-$j$ window, replace it by
\begin{equation}\label{eq:trace-joint-window}
 w_j(s)=w_j^0(s)+
   \sum_{d\in K_{j+1}}\sum_{i=0}^2
           h_{L_j,b_{d,i}+1}(x(3s)).
\end{equation}
The sum is finite at fixed $j$, so $w_j(s)=o(s)$. No uniform bound in $j$
is required by the common scale construction. Each summand is bounded
by the whole nonnegative window, which therefore pays every active
panel separately. The construction preserves the previous minimum
physical length and all previous windows and entry thresholds.

For completeness, this window also protects a chart whose level has
advanced beyond the paid stage-$j$ interval. Take $i_0=k_{j-j_*}$, the center of the constant-stage interval
\eqref{eq:common-long-plateau}, chosen with the enlarged entry threshold
and window. Thus $[i_0,i_0+w_j(s_{i_0})]$ has stage $j$.
Write $s=s_{i_0}$, and let $k$ be a later level with $J(k)=j+1$
and $t=s_k$. Put $H=h_{L_j,b}(x(t))$, where $b=b_{d,i}+1$.
Increase the stage-$j$ entry threshold beforehand so that, for every
later $t$,
\[
 h_{L_j,b}(x(t))\log L_{j+1}\le t/2,\qquad
 s_Gh_{L_j,b}(x(t))\log L_{j+1}\le x(t)/4.
\]
These inequalities are possible by the sublinear bound just proved.
If $t\le3s$, monotonicity of
$h_{L_j,b}$ on its large-$x$ domain gives
$H\le h_{L_j,b}(x(3s))\le w_j(s)$, and the end of the paid interval
lies strictly before $k$. Hence $H\le k-i_0$. If $t>3s$, use the first of these inequalities. Since block lengths are nondecreasing and the stages
from $i_0$ to $k$ are $j$ or $j+1$,
$t-s\le(k-i_0)\log L_{j+1}$. Now $t-s>2t/3>t/2$, so again
$H\le k-i_0$. Thus the backward depth paid by $H$ cannot leave the two
stages. Write $d_1,d_2$ for its first and second summands. All block
lengths between $i_0$ and $k$ are at least $L_j$. Hence, for every
reduced arc $0\le u\le b$,
\[
 \frac{a_{k-d_1}}{a_k}\ge L_j^{d_1}\ge\mathcal M(x(t))^{u+2}.
\]
The middle level therefore satisfies $\ell\ge k-d_1$. Since
$k-i_0\ge d_1+d_2$, the level $\ell-d_2$ is still at least $i_0$,
and
\[
 \frac{a_{\ell-d_2}}{a_\ell}\ge L_j^{d_2}\ge C(x(t))^{10}.
\]
The root depth consequently satisfies $w\ge\ell-d_2\ge k-H$.
This proves the dilation, depth and two-stage bracket. Since every block
length on $[w,k]$ is at most $L_{j+1}$,
\[
 \widehat\beta_k-\widehat\beta_w
   =s_G\log(a_w/a_k)\le s_GH\log L_{j+1}\le x(t)/4.
\]
Thus the second displayed entry inequality bounds the coupling decrease
in both cases. This proves the
$3/4$ plateau required in local input 2, not merely the weaker $1/2$
plateau used elsewhere. Include also the eventual monitor bound
$\max(L_j,L_{j+1})\le\mathcal M(x(t))$ in the entry threshold.

\emph{Offsets and common readings.}
Fix $d$. Since the stage map diverges, $d\in K_{J(k)}$ for all large
$k$. The preceding construction and local input 2 therefore give the
legacy costs for all three panels, uniformly for every smaller
nonnegative arc. Replace their arcs by $b_{d,i}+t$, $0\le t\le1$.
Their lower bounds are preserved because $b_{d,i}\ge d+6$; their other
data are unchanged. Intersect the finitely many eventual level sets
for the three panels with the set supplied by local input 3, and with
$k\ge\max_{e\le d}\operatorname{ref}(e)$. The resulting tail is
independent of $t$. Use the reading of panel $0$ for the chart maps.
For another panel retain exactly the same raw polynomial, Wick
polynomial and reading. The degree and coefficient bounds are unchanged.
The halo is unchanged because $Q_{d,j}$ and the budget radius agree;
the locality and healthy-set estimate are unchanged because the healthy
threshold also agrees. Finally the covariance and paid label norm belong
to the same history. Thus every clause of the stamp still holds, now
for this panel's representation. Its individual actuation witness and
this transported stamp certify the same chart. This is a reuse of one
stamp, not an assertion that independently chosen stamps have equal
readings. This proves
simultaneous availability, rather than three unrelated chart choices.
Every choice preceded the null sequence; the later interpolation and
subsequence selection only query these already constructed panels.
\end{proof}

That theorem also requires a simple gauge group, the funded chart-control
estimate, a bound on the generator through grade $d$, an index $l$ with
$l\leq\operatorname{trays}(d)$, a positive sequence tending to zero, and
the selected extraction order, smooth table, point limits, and reference
hierarchy used by the station construction.
It returns a strictly increasing level curve and, at a selected point,
data $(k,i,b,c,R)$ such that the requested scale equals the exact clock
gap at level $i$ and base arc $b$. Its chart has $5<b$ and $b$ no larger
than the permitted angular bound, covers all reference grades through
$d$, and supplies the finite-volume charts and invoices whenever $i<N$.
It also supplies a pairwise comparison receipt at cost
$\operatorname{costHint}(l)$ and a station whose lift has a finite norm
bound, preserves the unit, real structure, and signed action, and obeys
the short-test estimate for arity at most $l$, tests in the stated window,
and error $\operatorname{costHint}(l)$. The proof of this quantitative
upgrade is the separate three-panel chart result; this section uses that
exact output and proves the scalar gap interpolation and recursion. In
the gallery construction one takes $l=d$, so the chart cutoff, arity cap,
and finite grade cutoff agree. The clock gives, for each fixed base arc
$b\geq0$, a positive gap
\begin{equation}\label{eq:ex-lattice-firesky-gap}
 g_b(k)=a_k m_k^b,
\end{equation}
where $a_k>0$, $m_k\geq1$, $g_b(k)\to0$, and eventually
$g_b(k)\leq g_{b+1}(k+1)$. These scale-gap facts are the input to the
selection argument.

\begin{lemma}[Clock-gap interpolation]
\label{lem:ex-lattice-firesky-gap-interpolation}
Fix $b\geq0$. Suppose $g_b(k)>0$, $g_b(k)\to0$, and
$g_b(k)\leq g_{b+1}(k+1)$ for all $k\geq k_0$. For every $J\geq k_0$
there is $r>0$ such that for each $0<s\leq r$ there are $k>J$ and
$t\in[0,1]$ with $s=g_{b+t}(k)$.
\end{lemma}
\begin{proof}
Take $r=g_b(J)>0$. Given $0<s\leq r$, choose the least $k>J$ with
$g_b(k)\leq s$; it exists by convergence to zero. Minimality gives
$s\leq g_b(k-1)$, and the one-step comparison at $k-1$ gives
\[
 g_b(k)\leq s\leq g_b(k-1)\leq g_{b+1}(k).
\]
The function $t\mapsto g_{b+t}(k)=a_km_k^{b+t}$ is continuous on
$[0,1]$. The intermediate value theorem supplies the desired $t$.
\end{proof}

\subsection{The fields of a chart witness}
\label{subsec:ex-lattice-firesky-witness}

The producers of the pairwise receipt use several fields of the
finite-volume chart witness. We list them here, with the two scale
definitions they require. We then prove three lemmas: the corrected labels
solve their reading equation, the insertions lie inside the home chamber
at the root depth, and the clock of
Section~\ref{sec:common-scale-selection} supplies the dilation, depth,
bracket and plateau. The lawful history (F3) is supplied by
Corollary~\ref{cor:rank-trace-all-levels}, and the readings and stamps
of (F5) by Proposition~\ref{prop:rank-selected-tested-readout}, under
their stated funded construction hypotheses. In this subsection, $k$
denotes the level of the chart and $\mu_k$ the monitor $m_k$ of
\eqref{eq:ex-lattice-firesky-gap}. $J$ is the stage clock and $L_j$ are
the block lengths of Section~\ref{sec:common-scale-selection}, and
$c_k=\lceil(\log\widehat\beta_k)^2\rceil$ is the cover size of the clock
at level $k$. The monitor is $\mu_k=\mathcal M_k$ of
Section~\ref{sec:rank-root-normalization}, and $\widehat\beta_j$ is the
clock coupling of Section~\ref{sec:common-scale-selection}.

\paragraph{Two levels attached to a chart.}
Let the chart have base arc $b$, and put $m=\mu_k^b$. Its \emph{middle
level} and \emph{root depth} are
\[
 \ell=\max\Bigl\{j\le k:\frac{a_j}{a_k}\ge m\mu_k^{2}\Bigr\},\qquad
 w=\max\Bigl\{j\le\ell:\frac{a_j}{a_\ell}\ge c_k^{10}\Bigr\},
\]
with the convention $\max\varnothing=0$. Thus $w\le\ell\le k$. The middle
cell of the chart is a cell of spacing $a_\ell$. The root tests and the
chamber condition of the chart are taken on the centered tiling at depth
$w$, whose cells are indexed by the sites of the level-$w$ torus.

\paragraph{The witness fields.}
A chart witness at level $k$ with base arc $b$ supplies the following
fields. Here $M$ is the marker allowance of the chart, $N>k$ and $p$ are
arbitrary.
\begin{enumerate}
\item[(F1)] \emph{Allowance.} Four conditions hold:
 \begin{itemize}
 \item dilation: $m\mu_k^2\le a_\ell/a_k$, so the set defining $\ell$
  is nonempty;
 \item depth: $c_k^{10}\le a_w/a_\ell$;
 \item wheel: $\min(L_{J(k)},L_{J(k)-1})\ge2$, where the clock factor of a
  stage is its block length;
 \item bracket: $J(k)-1\le J(j)$ for every $w\le j\le k$.
 \end{itemize}
\item[(F2)] \emph{Plateau.} $\widehat\beta_k/2\le\widehat\beta_j\le\widehat\beta_k$
 for every $w\le j\le k$.
\item[(F3)] \emph{Lawful history.} For each $N>k$ and $p$ there are
 representations $P_j$ on the level-$j$ torus, $w\le j\le k$, along the
 course of the chart. Each $P_j$ satisfies the suppositions of
 Proposition~\ref{prop:rank-root-output-ticket}(2) at level $j$, with the
 clip and row bound of the stage $\pi_{J(j)}$. Its budget radius is
 $\lceil(\log\widehat\beta_j)^2\rceil^2$, and its threshold is
 \[
  \theta_{P_j}=\widehat\beta_j^{-1/2}\exp\bigl((\log\log\widehat\beta_j)^{3/2}\bigr),
 \]
 the clock's hatted coupling $\widehat\beta_j^{-1/2}$ at $j$ times the
 healthy allowance $\exp((\log\log\widehat\beta_j)^{3/2})$. This is the
 threshold of the representation itself. Once $\log\log\widehat\beta_j\ge1$ it is at most the clock
 threshold $\theta_{\rm clock}(j)=\widehat\beta_j^{-1/2}\mathcal M_j$ of
 Subsection~\ref{subsec:rank-root-clock-threshold}, in which the healthy
 allowance is replaced by the monitor $\mathcal M_j$; the numerical
 late-scale bounds there concern $\theta_{\rm clock}$.
\item[(F4)] \emph{Ticket.} $T_M(\pi_{J(w)})\le\log\mu_k/(16M)$.
\item[(F5)] \emph{Readings and stamps.} The level satisfies
 $\operatorname{ref}(d')\le k$ for every grade $d'\le d$, where
 $\operatorname{ref}$ is the reference level of
 Section~\ref{sec:actual-rank-purchases}. This holds when the chart is an
 output of the one-step construction, which covers all reference grades
 through $d$, and in any case for all large $k$, since only finitely many
 grades occur. For every grade $d'\le d$ there is a linear reading $T_{d'}$
 of the grade-$d'$ label space with three properties:
 \begin{itemize}
 \item it fixes the reference unit;
 \item it commutes with the real structure;
 \item it satisfies the half-accuracy bound
  $|a(v)-T_{d'}v|_1\le\frac12|a(v)|_1$, where $a=A_{k,N,p}$ is the
  calibrated entry (primary-coordinate) map at level $k$ for $(N,p)$.
 \end{itemize}
 It also carries a stamp duty with the same reading: a bound, in the
 halo of the readers, on the price of the insertions whose low-grade
 reading is $T_{d'}$. This clause is supplied for the funded calibrated family by
 Proposition~\ref{prop:rank-selected-tested-readout}; its hypotheses
 concern the same raw polynomial and reading, not separate choices.
\end{enumerate}
By (F1), every level of $[w,k]$ lies in the stage $J(k)$ or $J(k)-1$. Hence
$J(w)\ge J(k)-1$, and $w\to\infty$ as $k\to\infty$ because $J$ diverges.

\paragraph{Corrected labels.}
Write $\sigma=\sigma_{k,N,p}=(\beta_{N,p,k}^{1/2}\,\theta_k)^{-1}>0$ for
the shrink factor, where $\beta_{N,p,k}$ is the actual coupling of
Section~\ref{sec:actual-rank-purchases} and $\theta_k=\theta_{\rm clock}(k)$
is the clock threshold. Write $W_\rho$ for the linear map multiplying each
coordinate of a label by $\rho$ to its position degree, and $S_\sigma$ for
multiplication by $\sigma$ to the factor degree. Both act on the
grade-$d'$ label space $V_{d'}$. By Lemma~\ref{lem:ex-lattice-firesky-solves}
below, $T_{d'}$ is invertible. The \emph{corrected label} of a basis
vector $e$ of $V_{d'}$ is
\[
 X_{N,p}e=T_{d'}^{-1}\bigl(S_{\sigma}^{-1}W_me\bigr).
\]
These are the linear maps $X^{(i)}_{N,p}$ of input (G2) in
Subsection~\ref{subsec:trace-rounded-comparison}, for the chart selected at index
$i$.

\begin{lemma}[Corrected labels solve the reading equation]
\label{lem:ex-lattice-firesky-solves}
The reading $T_{d'}$ is bijective, and $X_{N,p}e$ is the unique label $v$
with
\begin{equation}\label{eq:ex-lattice-firesky-solves}
 S_\sigma\bigl(T_{d'}v\bigr)=W_me .
\end{equation}
\end{lemma}
\begin{proof}
Since $\operatorname{ref}(d')\le k<N$ by (F5), the calibrated entry map
satisfies the reference inverse estimate
\eqref{eq:rank-reference-inverse} of
Section~\ref{sec:actual-rank-purchases}, $|v|_1\le D_k|a(v)|_1$. With the
half-accuracy bound of (F5), the argument given there shows that $T_{d'}$
is injective, hence bijective. Next, $S_\sigma$ is invertible. So the
equation has the unique solution $X_{N,p}e$; this is
\eqref{eq:rank-common-lifts} with $\rho=m$.
\end{proof}

\paragraph{Status.} Fields (F1)--(F5) are fields of the funded witness.
For the clock of Section~\ref{sec:common-scale-selection} and all
sufficiently large $k$, with the requests of
\eqref{eq:rank-root-ticket-requests} included, (F1), (F2) and (F4) are
proved below: by Lemma~\ref{lem:ex-lattice-firesky-clock-fields}, the
remark after it, and Proposition~\ref{prop:rank-root-output-ticket}. The
inputs listed elsewhere that these fields supply are:
\begin{itemize}
\item the allowance and the root-depth identification used in
 Subsection~\ref{subsec:trace-rounded-receipt};
\item the bracket property used after
 Proposition~\ref{prop:rank-root-output-ticket};
\item the ticket at the output level.
\end{itemize}
Corollary~\ref{cor:rank-trace-all-levels} supplies (F3) on the entire
window, while Proposition~\ref{prop:rank-selected-tested-readout}
supplies (F5), under their stated funded hypotheses. These local results
do not replace the common all-grade purchase proved above.

\begin{lemma}[Chamber interiority at the root depth]
\label{lem:ex-lattice-firesky-chamber}
Consider a chart at level $k$ with base arc $b$, middle level $\ell$ and
root depth $w$. Let $x$ be fit for the chart: the grid configuration
$z=a_k\lceil mx\rceil$ is coordinatewise separated and
$|z_{c,\alpha}|\le a_\ell/8$ for $c=0,1$ and $\alpha=1,\dots,4$. Fix a grade
$d$, and let $Z_d$ bound the coordinates of the finite stencil $S_d$ of
the grade-$d$ insertions. Then, for every $N\ge k$ with
\begin{equation}\label{eq:ex-lattice-firesky-chamber-mesh}
 a_N(Z_d+2)<\tfrac38\,a_w ,
\end{equation}
every period $p$ and every grade-$d$ label, the insertion at the site
$z_c$ ($c=0,1$) has every link of its stencil, including both endpoints,
strictly inside the home chamber of the centered tiling at depth $w$. It
is therefore in the chamber-local domain of that tiling. Condition
\eqref{eq:ex-lattice-firesky-chamber-mesh} holds for all large $N$, with
a threshold depending only on $w$ and $d$. In particular (G4) of
Subsection~\ref{subsec:trace-rounded-comparison} holds with $w_i=w$ for
the chart selected at index $i$.
\end{lemma}
\begin{proof}
Measure positions in mesh steps of the level-$N$ lattice; since $N\ge k$,
$a_k/a_N$ is an integer. The depth-$w$
cell has an odd number $q=a_w/a_N$ of steps along each axis, as in the
interiority argument of Section~\ref{ex-centered-construction}. Its home
cell consists of the offsets $u$ with $|u_\alpha|\le(q-1)/2$ from the
centre. Hence the site $z_c$ has integer
offset $u=z_c/a_N$. By fitness, and because $w\le\ell$ implies
$a_\ell\le a_w$,
\[
 |u_\alpha|\le\frac{a_\ell}{8a_N}\le\frac{q}{8}.
\]
A link of the stencil at $z_c$ has endpoints $u+s$ and $u+s+e$, with
$s\in S_d$ and $e$ a unit vector. Each coordinate of an endpoint is at
most $q/8+Z_d+1$ in absolute value. By
\eqref{eq:ex-lattice-firesky-chamber-mesh},
$Z_d+2<3q/8$. Hence
\[
 \frac q8+Z_d+1<\frac q2-1<\frac{q-1}2 ,
\]
so both endpoints lie strictly inside the home cell. The insertion is a
bounded, measurable, gauge-invariant function of these links. It is
therefore chamber-admissible, exactly as in the argument cited above.
Finally, $a_N\to0$ while $a_w>0$ and $Z_d$ are fixed, so
\eqref{eq:ex-lattice-firesky-chamber-mesh} holds for all large $N$. The
threshold does not depend on $p$, on the label, or on the fit
configuration.
\end{proof}

\begin{lemma}[The clock supplies the dilation, depth, bracket and plateau]
\label{lem:ex-lattice-firesky-clock-fields}
Let the clock be the one of Section~\ref{sec:common-scale-selection}, so
that $\widehat\beta_m=\beta_*+s_Gs_m$, $a_m=e^{-s_m}$,
$s_{m+1}=s_m+\log L_{J(m+1)}$, $L_j\ge3$, and $J=j_*+J_0$. Let the
monitor be $\mu_k=\exp((\log\log\widehat\beta_k)^2)$ and the cover size
$c_k=\lceil(\log\widehat\beta_k)^2\rceil$. Fix $\bar b>0$. Then for all
sufficiently large $k$ and every base arc $0<b\le\bar b$, the following
hold.
\begin{enumerate}
\item The sets defining $\ell$ and $w$ are nonempty. Hence
 $m\mu_k^2\le a_\ell/a_k$ and $c_k^{10}\le a_w/a_\ell$; these are the
 dilation and depth of (F1).
\item $J(k)-1\le J(j)$ for every $w\le j\le k$; this is the bracket of
 (F1).
\item $\widehat\beta_k/2\le\widehat\beta_j\le\widehat\beta_k$ for every
 $w\le j\le k$; this is (F2).
\end{enumerate}
\end{lemma}
\begin{proof}
\emph{Growth of the scale.} At every step $m$ with $J(m+1)>j_*$, the
entry requirement \eqref{eq:common-stage-entries} gives
$\log L_{J(m+1)}\le s_m/J(m+1)$, so $s_{m+1}\le2s_m$ at each such step.
The steps of the first plateau are finitely many. Hence $s_k\le C_02^k$ for a constant $C_0$, and
\[
 \log\log\widehat\beta_k\le\log k+C_1
\]
for large $k$, with a constant $C_1$. On the other side,
$s_k\ge k\log3$ by \eqref{eq:common-scale-brackets}.

\emph{Item 1.} The set defining $\ell$ contains $j=0$ as soon as
$a_0/a_k=e^{s_k}\ge m\mu_k^2=\mu_k^{b+2}$. Now
$e^{s_k}\ge3^k$, while
$\mu_k^{b+2}\le\exp\{(\bar b+2)(\log k+C_1)^2\}$, so this holds for large
$k$. The defining inequality then holds at its maximum $\ell$. If $\ell<k$, maximality gives
$3^{k-\ell-1}\le a_{\ell+1}/a_k<\mu_k^{b+2}$, because
$a_{\ell+1}/a_k=\prod_{t=\ell+2}^kL_{J(t)}$ and $L_j\ge3$. Likewise, if
$w<\ell$, then $3^{\ell-w-1}\le a_{w+1}/a_\ell<c_k^{10}$, once $w$ is
defined. Hence
\[
 k-\ell\le1+\frac{(\bar b+2)(\log k+C_1)^2}{\log3},\qquad
 \ell-w\le1+\frac{10\log c_k}{\log3}.
\]
The first bound gives $k-\ell=O((\log k)^2)$, so $\ell/k\to1$.
Since $c_k\le2(\log\widehat\beta_k)^2$ for large $k$ and
$\log\log\widehat\beta_k\le\log k+C_1$, we have
$\log c_k=O(\log k)$. Thus $10\log c_k=o(\ell)$, and for large $k$,
$a_0/a_\ell\ge3^\ell\ge c_k^{10}$. The set defining $w$ is therefore
nonempty. Adding the two bounds gives a constant $C_2$ with
\[
 k-w\le C_2(\log k)^2 .
\]

\emph{Item 2.} Let $j=J_0(k)$, so $A_j\le k<A_{j+1}$ in the notation of
Section~\ref{sec:common-scale-selection}. If $w\ge A_j$, every level of
$[w,k]$ has $J_0=j$. Otherwise $w<A_j\le k$, and so
$k-A_j\le k-w\le C_2(\log k)^2$. For large $k$ this forces
$k\le2A_j$. By the recursion,
\[
 A_j-A_{j-1}=b_j+2k_{j-1}+1\ge4A_{j-1}+3 ,
\]
so $A_j-A_{j-1}\ge\tfrac45A_j$. For large $k$, and hence large $A_j$,
\[
 k-w\le C_2(\log2A_j)^2\le\tfrac45A_j\le A_j-A_{j-1}.
\]
Therefore $w\ge k-(A_j-A_{j-1})\ge A_{j-1}$, and every level of $[w,k]$
has $J_0\ge j-1$.

\emph{Item 3.} Since $s_m$ is increasing, $\widehat\beta_j\le\widehat\beta_k$
for $j\le k$, and $\widehat\beta_j\ge\widehat\beta_w$ for $j\ge w$. Next,
$\widehat\beta_k-\widehat\beta_w=s_G\log(a_w/a_k)$. Split $a_w/a_k$ into
four factors:
\[
 \frac{a_w}{a_k}
 =\frac{a_w}{a_{w+1}}\cdot\frac{a_{w+1}}{a_\ell}\cdot
  \frac{a_\ell}{a_{\ell+1}}\cdot\frac{a_{\ell+1}}{a_k}.
\]
Omit the factors whose indices coincide. By maximality the second and
fourth factors are below $c_k^{10}$ and $\mu_k^{b+2}$. The first and third
are block lengths $L_{J(t+1)}$ with $t\ge w$. By the entry requirement
each is at most $\exp(s_k/J(t+1))\le\exp(s_k/(J(k)-1))$, using item 2.
Hence
\[
 \log\frac{a_w}{a_k}\le\frac{2s_k}{J(k)-1}+10\log c_k
   +(\bar b+2)(\log k+C_1)^2 .
\]
Since $J(k)\to\infty$ and $s_k\ge k\log3$, the right side is at most
$s_k/2$ for large $k$. Thus
$\widehat\beta_k-\widehat\beta_w\le s_Gs_k/2\le\widehat\beta_k/2$,
because $\widehat\beta_k\ge s_Gs_k$.
\end{proof}

The wheel bound of (F1) is also immediate: the clock factor of a stage
$j$ is its block length $L_j\ge3$. Consequently, for all sufficiently
large $k$ and base arcs up to $\bar b$, of the witness fields,
(F3) and (F5) are supplied by the conditional constructions just cited;
the remaining issue is to pay their hypotheses on the same selected
history, as in Proposition~\ref{prop:trace-joint-chart-purchase}. (F4) is replaced by
Proposition~\ref{prop:rank-root-output-ticket}.

\subsection{Choosing scales and chart levels}

\begin{lemma}[Strict scale selection]
\label{lem:ex-lattice-firesky-selection}
Assume the one-step chart producer has this output at every grade $d$: for
every previous index $N$ and every $\theta>0$, it supplies an index $n>N$,
a level $i\geq d$, and a grade-$d$ chart with exact radius $s_n<\theta$
and the finite-level receipts above. Let $s_n>0$ and $s_n\to0$. Then
there are strictly increasing indices $n_d$, levels $i_d\geq d$, and
admissible charts at grade $d$ and level $i_d$, such that
\[
 s_{n_0}<1,\qquad s_{n_{d+1}}<s_{n_d},\qquad i_d\geq d.
\]
In particular $i_d\to\infty$ and the selected chart radius is exactly
$s_{n_d}$.
\end{lemma}
\begin{proof}
For $d=0$, apply the one-step output with previous index $0$ and threshold
$1$. For $d>0$, apply it with previous index $n_{d-1}$ and threshold
$s_{n_{d-1}}$. It gives $n_0>0$, $s_{n_0}<1$, and thereafter
$n_d>n_{d-1}$, $s_{n_d}<s_{n_{d-1}}$, with a grade-$d$ chart at level
$i_d\geq d$ and exact radius $s_{n_d}$. Induction proves the claims;
$i_d\geq d$ implies $i_d\to\infty$.
\end{proof}

The one-step chart theorem obtains this output by combining the clock-gap
interpolation with the funded chart family. Its quantitative comparison
and station estimates are the three-panel weak-comparison result in
Section~\ref{sec:three-panel-weak-comparison}; we use its conclusion here
without duplicating those estimates.

Here is the finite-level extraction used to pass from the charts to the
limiting fields. Let $F$ be an extraction order, let $\mathcal U$ be an
ultrafilter on period-phase pairs finer than the fused filter of $F$, and
fix a cutoff $D$. Suppose $X_{N,p}$ is a tuple of linear maps, one for
each grade $0,\ldots,D$, from the corresponding reference field into the
common finite-period coefficient space. Assume $A\geq0$ and, on an
$F$-late set,
\[
 \|X_{N,p}\|\leq A.
 \tag{1}\label{eq:ex-lattice-firesky-lift-bound}
\]
The finite product of these map spaces is finite-dimensional. The closed
ball of radius $A$ is compact. Compactness gives a proper subfilter limit
of the bounded tuple of maps. This subfilter is finer than $\mathcal U$;
ultrafilter maximality identifies it with $\mathcal U$. Thus the tuple
converges along $\mathcal U$ to $Y=(Y_0,\ldots,Y_D)$. This is the only
compactness step: it is finite-dimensional and is performed at each fixed
cutoff.

If the finite-level maps preserve the real involution and signed actions on
the late set, continuity of those actions shows that each $Y_d$ preserves
them. If they send the reference unit to the scalar unit eventually, the
same limit identity gives $Y_d(1)=1$. Compose with the quotient projection
to define the table map
\[
 U_d=\mathsf P_dY_d,\qquad 0\leq d\leq D.
\]
This order matters: the finite maps converge in the coefficient space, and
the table-valued maps are obtained by applying the continuous projection
after taking the limit.

For a fixed tuple of $n$ labels of grades at most $D$, a fixed smooth test
function supported in the allowed window, and $n$ no greater than the
available arity cap, the finite-level receipt bounds the paid test error
by $\epsilon$ times the label charge and the test supremum. To pass this
estimate to the quotient table, use the ordered-moment limit rather than
continuity of an integral whose finite-volume measure varies. For each
fixed tuple of coefficient vectors, the ordered read receipt gives
convergence of the corresponding finite-volume moment to the table moment.
The moment is multilinear in its $n$ label vectors, and each label space is
finite-dimensional. Expanding each moving vector in a fixed basis expresses
the moment as a finite sum of products of convergent coordinates and
convergent fixed-basis moments. Hence the moving-label moment converges to
the table moment. This is the finite-basis multilinear limit argument; it
uses the common ordered extraction and the fixed compactly supported test.
The inequality passes to the limit
because a closed interval is closed under limits. More explicitly, for
each $\eta>0$ the late-set estimate is at most the claimed bound plus
$\eta$; taking the limit and then letting $\eta\downarrow0$ removes the
slack. This proves the cutoff-level short-test estimate with no exchange
of an infinite sum and a limit. It is the same moment-passage mechanism as
the calibrated-smearing lemma: fixed-vector convergence is supplied by the
ordered table, and finite dimensionality handles the moving vectors.

Apply this extraction separately at cutoff $D=d$ to the chart chosen at
grade $d$ and level $i_d$. The chart receipt supplies a bound on its finite lifts,
exact reality and signed covariance, the exact unit, and a short-test error
for arity $n\leq d$ with coefficient $1/(d+1)$. Projecting the
limits gives panels $U_d$. For every fixed finite label tuple, its grades
are at most $d$ eventually along the sequence of cutoffs; the support of a
fixed compact test fits the expanding chart window eventually; and
$1/(d+1)\to0$. Write
$\mathcal T_{r,U}(v;t)$ for the tested integral of the tuple $v$ under
radius $r$, panel family $U$, and test function $t$. Write
$\mathcal T_0(v;t)$ for the corresponding reference test value. Hence
\[
 \mathcal T_{s_{n_d},U_d}(v;t)-\mathcal T_{0}(v;t)
 \longrightarrow0
\]
for every fixed tuple $v$ and test $t$. This is the tested approach property.
The argument is for each fixed test; it does not assert a uniform tail over
all tests or all label tuples.

The final chart isomorphisms follow grade by grade. For a fixed grade $d$,
the reference tests separate its finite-dimensional field: every nonzero
vector has a nonzero separated two-point evaluation. To obtain a smooth
test, keep the spectator insertion fixed and choose a separated pair where
the complex value is nonzero. Multiply the spectator by a unit complex
phase so that the real part of the value is positive. Continuity away from
the diagonal gives a small neighborhood on which that real part remains
positive. Integrating against a nonnegative smooth bump of integral one,
supported in that neighborhood, gives a smooth smeared test with nonzero
value. Thus smooth smeared reference tests also separate. Normalize to the
unit sphere. For
each point of that compact sphere, continuity gives a neighborhood on
which one such test has absolute value at least half its value at the
point. A finite subcover gives tests $t_1,\ldots,t_m$ and a constant
$\alpha>0$ such that every unit vector has
$\max_j|\mathcal T_0(v;t_j)|\geq\alpha$. The fixed finite set of
label charges and test suprema is bounded on that sphere. The tested
approach estimate therefore makes all corresponding errors less than
$\alpha/2$ for sufficiently large $k$. If the projected grade map killed a
unit vector, all its tested values would vanish, contradicting this lower
bound. Thus it is injective. The dimension bound from the selected table gives the reverse
dimension inequality; therefore the two dimensions agree and the map is
bijective. This tail may depend on $d$: there is no single finite index
after which every grade is simultaneously bijective. For the finitely many
indices before the grade-$d$ tail, choose any isomorphism using the same
dimension equality. This defines a scaling at every index. For each fixed
grade, the unit, reality, and signed symmetry identities hold on its
eventual isomorphism tail; the arbitrary isomorphisms used earlier need not
preserve those identities. The reference test-separation input here is the
all-grade free-detector result of
Section~\ref{sec:ex-lattice-all-grade-free-detectors}; the separate
low-grade trace calculation is not enough to justify separation in every
grade. The reviewed corrected-lift argument in
Section~\ref{sec:ex-lattice-trace-producer}
provides the compatible panel maps when required by the trace family.

\subsection{The gallery handoff}

The preceding argument proves the scale-selection, finite-cutoff extraction,
tested-approach, and fixed-grade eventual-isomorphism clauses. It does not
by itself prove the complete gallery statement. That handoff also needs the
physical-array receipts in the selected-run transfer. For the same
ultrafilter, hierarchy, and scaling these include the following. Let
$b_i>0$ be the selected arch scale at the $i$th postselected level, so
that $b_i\to\infty$; let $1$ and $\mathcal O$ be the unit and the
normalized curvature-square field of the reference grade-four table, and
$d_2$ the physical coefficient, as in
Section~\ref{sec:ex-lattice-trace-scribe-producer}. For a comparison table
$T$, write $\Theta_T$ for the scalar trace of the physical stress centered
by its one-point value in $T$, and $U_{i,4}$ for the grade-four chart from
the reference table to the selected table on its eventual isomorphism
tail. The receipts are:
convergence of the recentered physical stress to its prescribed principal
material tensor; for every comparison table that obeys the table hierarchy
laws and respects the selected table, real sequences $d_i,e_i$ such that,
eventually in $i$, the centered scalar trace satisfies
\[
 r_i^4U_{i,4}^{-1}\Theta_T=e_i\,1+d_i\,(4\mathcal O),
 \qquad b_id_i\longrightarrow d_2,\qquad b_ie_i\longrightarrow0.
\]
Thus $d_i,e_i\to0$ themselves. The final receipt is, for every comparison
table, convergence of $b_ir_i^4U_{i,4}^{-1}\Theta_T$ to
$4d_2\,\mathcal O$. The second direction
$4\mathcal O$ is the trace of the true-branch material tensor; the limit
of the zoomed observable is a multiple of $\mathcal O$, not of the unit. The arbitrary-null-
sequence case needs these same receipts after the selected scaling has
been constructed above. They are additional physical and field-array
estimates, not consequences of finite-dimensional compactness or of the
rank-one solve alone. The correction, one-dimensional
kernel, stable inverse, flux normalization, and physical/reference shell
comparison in Section~\ref{sec:ex-lattice-trace-producer} remain conditional
on their displayed analytic receipts.

In the selected public construction, the one-scale chart family comes from
the paid multi-scale request and its subsequent chart staging. This section
proves the scale-selection and tested-chart clauses of that assembly; the
physical-array, stress, and flux receipts needed for the full gallery remain
separate inputs.
